\documentclass[UTF8,12pt,fontset=none]{ctexart}
\usepackage[a4paper,margin=1in]{geometry}
\usepackage{amsmath}
\usepackage{array}
\usepackage{booktabs}
\usepackage{enumitem}
\usepackage{hyperref}
\usepackage{longtable}
\usepackage{xcolor}

\setCJKmainfont{Songti SC}
\setCJKsansfont{PingFang SC}
\setCJKmonofont{STHeiti}

\hypersetup{
  colorlinks=true,
  linkcolor=black,
  citecolor=blue,
  urlcolor=blue,
}

\setlist[itemize]{leftmargin=2em}
\setlength{\parskip}{0.4em}

\title{看见什么、重建什么、想象什么:\\
面向视觉语言模型时空知觉的行为反演研究提纲}
\author{ordinary-bench 论文规划稿}
\date{2026年3月12日}

\begin{document}
\maketitle
\tableofcontents

\section{推荐定位与边界}

\paragraph{推荐主标题}
\textbf{看见什么、重建什么、想象什么：面向视觉语言模型时空知觉的行为反演研究}

\paragraph{一句话摘要}
本文不把视觉语言模型（VLM）当作黑箱答题器来比较准确率，而是借鉴神经科学和心理物理学中的\textbf{行为反演}思路，把模型回答转化为可检验的时空约束，再通过场景重建、场景补全和未来想象任务，分析模型究竟从图像和视频中\textbf{看见了什么}、\textbf{稳定保留了什么}、以及\textbf{无法形成何种时空表征}。

\paragraph{推荐主线}
第一版主文应聚焦三件事：
\begin{itemize}
  \item 一个建立在 3D 渲染世界上的\textbf{受控时空知觉评测框架}；
  \item 一个从 VLM 行为反推出\textbf{外显场景信念（scene belief）}的重建分析框架；
  \item 一套区分\textbf{真实视觉贡献}与\textbf{语言先验}的控制实验。
\end{itemize}

\paragraph{当前仓库的真实边界}
当前仓库最成熟的是静态空间基准、三条件控制实验、多视角输入与 2D scene belief reconstruction。动态部分已经具备\textbf{可生成短视频、逐帧状态与事件标签}的原型，但主链路尚未完全闭环。因此文稿中建议使用如下措辞：
\begin{itemize}
  \item 强调 \textbf{3D-rendered controlled worlds}，而不是宣称已经完成 full 3D world model benchmark；
  \item 强调 \textbf{2D/2.5D 的外显时空表征分析}，而不是宣称已经恢复真实三维物理状态；
  \item 把动态部分写成\textbf{时空扩展主实验或强补充实验}，不要写成“所有模块都已 fully scaled”。
\end{itemize}

\section{核心科学问题与研究假设}

\subsection{核心问题}
\begin{itemize}
  \item \textbf{RQ1:} VLM 的空间与时间回答中，有多少信息真正来自视觉输入，而不是语言先验或任务模板？
  \item \textbf{RQ2:} 当问题是人类直观且低门槛的时空判断时，VLM 是否能形成稳定、一致、可重建的场景表征？
  \item \textbf{RQ3:} VLM 的表征更像是局部 2D 纹理匹配，还是能够支持物体级、关系级、轨迹级的时空重构？
  \item \textbf{RQ4:} 在场景重绘、场景补全和未来想象任务中，VLM 暴露出的误差主要来自感知缺失、结构不一致，还是时序建模失败？
  \item \textbf{RQ5:} 多视角、长时间跨度、顺序打乱、时间反转等操控，会如何改变 VLM 的时空信念？
\end{itemize}

\subsection{研究假设}
\begin{itemize}
  \item \textbf{H1:} 单题准确率会高估 VLM 的几何能力；一旦转向场景级重建，模型的真实时空表征缺陷会被放大。
  \item \textbf{H2:} 允许拒答、集合回答与重述一致性检查后，模型的约束质量会优于强制作答协议\cite{whitehead2022reliablevqa,eisenschlos2024selectivevqa,khan2024consistency}.
  \item \textbf{H3:} 正确图像条件下的场景重建质量，应显著优于错误图像与无图像条件，否则说明模型主要依赖语言先验\cite{lee2025vlind,chen2024mmstar}.
  \item \textbf{H4:} 空间关系任务中，距离排序能力会明显优于方向、顺序、速度与轨迹理解；多视角能显著缓解静态歧义，但不能自动解决时间知觉问题\cite{chen2024spatialvlm,fu2024blink}.
  \item \textbf{H5:} 即使 VLM 能给出局部正确答案，它仍可能无法形成单模态、低 spread 的场景信念，表现为高冲突、低可满足性和多模态解空间。
\end{itemize}

\section{摘要提纲}

\subsection{问题动机}
现有 VLM 基准普遍将视觉理解压缩为离散问答分数，难以区分“真正看见了”“依赖语言先验猜对了”“知道自己看不清”的差异。与此同时，近期工作显示，许多大模型在视觉细粒度感知、语言先验剥离和视频时序理解上仍存在系统性盲点\cite{chen2024spatialvlm,fu2024blink,chen2024mmstar,lee2025vlind,fu2024videomme,zhou2024mlvu,liu2024tempcompass}.

\subsection{方法概述}
本文提出一个受神经科学行为实验启发的 VLM 时空知觉评测框架：在 3D 渲染环境中生成静态场景和动态短视频，构造空间关系、时间关系、场景重绘、场景补全和未来想象任务；再将模型输出转为关系约束，通过 scene belief reconstruction 反推出模型的外显时空表征，并用错误图像、无图像、视角变化、顺序打乱与时间反转等对照实验估计真实视觉信息增益。

\subsection{预期结论}
预期发现包括：VLM 对静态距离排序的表现显著优于方向与时序任务；多视角能改善空间感知但难以替代稳定的时间建模；在错误图像和无图像条件下，模型仍会输出具语言偏置的“伪场景信念”；场景重建指标比单题准确率更能揭示 VLM 的真实时空缺陷。

\subsection{意义}
本文的意义不在于给 VLM 再加一个 benchmark，而在于提出一种\textbf{从行为反推内部时空表征}的方法论，使“VLM 看到了什么”可以被操作化、量化并与生物启发的实验范式接轨。

\section{论文整体结构}

\subsection{第 1 节：引言}
引言建议按如下逻辑展开：
\begin{itemize}
  \item 先指出大模型评测的核心缺口：我们知道它们能答题，但不知道它们是否形成了稳定的时空世界表征。
  \item 再指出现有视觉基准的问题：任务往往偏向识别、语言化描述或粗粒度视频 QA，缺少对\textbf{空间几何、时序连续性、场景可重建性}的联合分析\cite{chen2024spatialvlm,fu2024blink,fu2024videomme,zhou2024mlvu,liu2024tempcompass,zhou2025vlm4d}.
  \item 接着提出本文视角：借鉴神经科学中的逆向研究思路，把模型行为视作观测信号，通过约束、重建和干预实验反演 VLM 的外显场景信念。
  \item 最后给出贡献列表。
\end{itemize}

\subsection{第 2 节：相关工作}

\subsubsection{VLM 的空间知觉与视觉盲点}
这一段建议覆盖：
\begin{itemize}
  \item 空间推理专门能力：SpatialVLM\cite{chen2024spatialvlm}；
  \item 视觉感知与推理脱节：BLINK\cite{fu2024blink}；
  \item 视觉评测路径本身是否合理：MMStar\cite{chen2024mmstar}；
  \item 语言先验污染视觉判断：VLInd-Bench\cite{lee2025vlind}。
\end{itemize}

\subsubsection{视频理解与时间推理}
这一段建议覆盖：
\begin{itemize}
  \item 通用视频评测：Video-MME\cite{fu2024videomme}；
  \item 长视频与多任务理解：MLVU\cite{zhou2024mlvu}；
  \item 专门面向时间推理的评测：TempCompass\cite{liu2024tempcompass}；
  \item 时间感知型模型：TimeChat\cite{ren2024timechat} 与 VTimeLLM\cite{huang2024vtimellm}；
  \item 更接近 4D 世界理解的视角：VLM4D\cite{zhou2025vlm4d} 与 Feature4X\cite{zhou2025feature4x}。
\end{itemize}

\subsubsection{选择性回答、可靠性与视觉贡献分离}
这一段建议强调：
\begin{itemize}
  \item 允许 abstain 的 VQA 可靠性\cite{whitehead2022reliablevqa}；
  \item selective VQA\cite{eisenschlos2024selectivevqa}；
  \item 黑盒 VLM 的一致性与不确定性分析\cite{khan2024consistency}。
\end{itemize}

\subsubsection{从约束到场景：重建与几何理论}
这一段建议作为方法学支撑：
\begin{itemize}
  \item ordinal embedding 与非度量嵌入\cite{agarwal2007gnmds,ma2018robustordinal}；
  \item 欧氏距离矩阵与可实现性\cite{dokmanic2015edm}；
  \item bearing / angle rigidity 对方向约束唯一性的解释\cite{zhao2019bearingrigidity,chen2021anglerigidity}。
\end{itemize}

\subsubsection{神经科学启发的模型分析}
这里不必把篇幅写得很重，但建议明确：
\begin{itemize}
  \item 借用“由行为观测反推内部表征”的思想，而不是声称能直接读出模型内部语义；
  \item 可用 \textit{Language Models Represent Space and Time}\cite{gurnee2024spacetime} 作为“模型内部可能形成时空结构”的背景；
  \item 若后续有开放权重模型，可在补充材料里加隐藏状态或表示几何分析；若没有，则主文坚持\textbf{行为可检验}路线。
\end{itemize}

\subsection{第 3 节：问题定义与评测对象}

\subsubsection{外显场景信念}
建议定义：
\[
B_m(I) \quad \text{或} \quad B_m(I_{1:T})
\]
表示模型 $m$ 对静态图像或视频输入的外显场景信念。这个对象不是隐藏向量，而是可由回答、拒答、重述一致性、集合式答案与场景重建共同反演出来的\textbf{可行约束结构}。

\subsubsection{三类能力的区分}
建议把论文中的能力拆成三类：
\begin{itemize}
  \item \textbf{感知（perception）}: 模型能否从图像或视频中读出局部空间与时间关系；
  \item \textbf{重建（reconstruction）}: 模型给出的约束能否支持稳定、低歧义的场景恢复；
  \item \textbf{想象（imagination）}: 在遮挡、缺帧、未来时刻预测等条件下，模型能否补全潜在场景。
\end{itemize}

\subsubsection{推荐写法}
这一节建议明确一句：
\begin{quote}
本文关注的不是 VLM 是否“会说关于空间和时间的话”，而是它是否形成了可由行为外显、可由重建检验、可由对照实验验证的时空世界模型。
\end{quote}

\subsection{第 4 节：数据、任务与协议设计}

\subsubsection{静态场景模块}
建议以当前仓库为主：
\begin{itemize}
  \item 在 Blender 中生成多物体场景、单视角与多视角图像；
  \item 使用 QRR（距离排序）和 TRR（相对钟面方向）刻画空间结构；
  \item 分析随物体数量上升的性能退化与可重建性变化。
\end{itemize}

\subsubsection{动态场景模块}
建议写成“动态扩展主实验”：
\begin{itemize}
  \item 从静态对象布局扩展到 static / linear / circular 运动轨迹；
  \item 逐帧记录位置、速度和事件标签；
  \item 在论文措辞上强调\textbf{由 3D 渲染环境生成的 2.5D 动态轨迹世界}，避免对 full 3D dynamics 过度承诺。
\end{itemize}

\subsubsection{四类任务族}
建议把任务分为四组：
\begin{itemize}
  \item \textbf{空间判断}: QRR / TRR / 多视角空间关系；
  \item \textbf{时间判断}: 起止关系、顺序、速度比较、轨迹类型、事件反转；
  \item \textbf{场景重绘与复原}: 输出俯视布局、对象关系图、遮挡补全；
  \item \textbf{未来想象}: 下一帧、终态、碰撞前后、关系翻转时刻预测。
\end{itemize}

\subsubsection{三种问答协议}
建议保留：
\begin{itemize}
  \item \textbf{forced QA}: 传统强制作答；
  \item \textbf{selective QA}: 允许拒答、模糊、集合回答\cite{whitehead2022reliablevqa,eisenschlos2024selectivevqa,khan2024consistency}；
  \item \textbf{direct extraction}: 直接输出对象关系集合或场景图，可参考 scene graph generation 路线\cite{kim2024llm4sgg,dutta2025openworldsgg}。
\end{itemize}

\subsubsection{关键控制实验}
建议至少写出以下干预：
\begin{itemize}
  \item 正确图像 / 错误图像 / 无图像；
  \item 单视角 / 多视角；
  \item 单帧 / 多帧；
  \item 原序视频 / 打乱帧序 / 时间反转；
  \item 完整场景 / 局部遮挡 / 缺帧补全。
\end{itemize}

\subsection{第 5 节：行为反演与神经科学式分析方法}

\subsubsection{从回答到约束}
建议说明：模型输出首先被转为关系观测，再归一化为 QRR、TRR、速度排序、事件顺序、轨迹类别等时空约束。

\subsubsection{从约束到场景重建}
建议把当前仓库中的 reconstruction 写成论文方法核心之一：
\begin{itemize}
  \item 距离约束用 ordinal embedding 视角建模\cite{agarwal2007gnmds,ma2018robustordinal}；
  \item 可实现性与几何自由度由 EDM 与 rigidity 理论解释\cite{dokmanic2015edm,zhao2019bearingrigidity,chen2021anglerigidity}；
  \item 输出场景级指标，而不是停留在单题分数。
\end{itemize}

\subsubsection{行为层面的“神经科学式”协议}
这里建议把方法论写成四个词：
\begin{itemize}
  \item \textbf{Probe}: 用重述、一致性、互逆、传递与跨条件对照，测试表征是否稳定；
  \item \textbf{Lesion}: 用遮挡、错图、删帧、乱序与反转，测试表征依赖于哪些视觉线索；
  \item \textbf{Reconstruct}: 用 scene belief reconstruction 测试这些行为能否支撑一个统一几何解释；
  \item \textbf{Imagine}: 用补全和未来预测测试模型是否拥有超出当前感知的时空内推能力。
\end{itemize}

\subsubsection{可选的内部表示分析}
如果后续只使用 API 模型，这一部分不要写成主方法，只可在讨论或补充材料中写成“未来可扩展”：
\begin{itemize}
  \item 开放权重模型可增加层级 probe、表征相似性或时间步激活分析；
  \item 主文结论应坚持建立在\textbf{可复现实验行为}之上。
\end{itemize}

\subsection{第 6 节：实验设计与指标体系}

\subsubsection{实验矩阵}
\begin{longtable}{p{0.16\linewidth}p{0.24\linewidth}p{0.28\linewidth}p{0.22\linewidth}}
\toprule
模块 & 研究问题 & 关键对照 & 核心指标 \\
\midrule
静态空间知觉 & 模型能否正确读取距离与方向结构 & 单视角 vs 多视角；小场景 vs 大场景 & QRR/TRR accuracy，missing rate \\
\midrule
视觉贡献分离 & 视觉信息是否真正在起作用 & 正确图像 vs 错误图像 vs 无图像 & VIG，vig\_tau，vision\_helps\_rate \\
\midrule
场景信念重建 & 回答能否支撑统一几何场景 & GT 约束 vs 模型约束；强制作答 vs selective QA & CSR，NRMS，Kendall $\tau$，$K_{\text{geom}}$，spread \\
\midrule
时间知觉 & 模型是否理解顺序、速度和轨迹 & 单帧 vs 多帧；原序 vs 乱序 vs 反转 & order accuracy，velocity ranking，trajectory accuracy \\
\midrule
场景补全与想象 & 模型是否能补全隐藏状态和未来状态 & 遮挡 vs 无遮挡；缺帧 vs 完整输入 & completion error，future event accuracy，trajectory NRMS \\
\bottomrule
\end{longtable}

\subsubsection{基础指标}
建议保留当前仓库已实现的指标：
\begin{itemize}
  \item 准确率：QRR、TRR、temporal relation、trajectory；
  \item 结构一致性：transitivity、reciprocity、跨重述一致性；
  \item 重建指标：CSR、NRMS、Kendall $\tau$、$K_{\text{geom}}$、spread；
  \item 视觉增益：正确图像相对错误图像和无图像的重建优势。
\end{itemize}

\subsubsection{人类基线}
建议单列一个小节：
\begin{itemize}
  \item 人类应在静态空间判断与粗粒度时间顺序上接近天花板；
  \item 若 VLM 在人类直观任务上仍显著失分，会强化论文的科学说服力；
  \item 人类基线应至少覆盖：静态空间、时间顺序、遮挡补全三类任务。
\end{itemize}

\subsection{第 7 节：结果组织建议}

\subsubsection{主结果}
主结果建议按“从简单到复杂”的顺序组织：
\begin{enumerate}
  \item 静态空间任务中，模型能做对哪些局部判断；
  \item 场景级 reconstruction 暴露出哪些全局不一致；
  \item 三条件实验显示多少能力是真正视觉驱动；
  \item 动态任务中，模型在哪些时间知觉维度急剧退化；
  \item 场景重绘、复原和想象任务如何进一步揭示表征边界。
\end{enumerate}

\subsubsection{预期结果图景}
建议预期如下图景：
\begin{itemize}
  \item QRR 明显高于 TRR，静态多视角能显著提高空间关系读取；
  \item 在错误图像和无图像条件下，模型仍会形成语言偏置驱动的伪约束；
  \item 从单题准确率看似“还不错”的模型，在 scene-level reconstruction 上可能不可行、欠约束或多模态；
  \item 动态任务中，长时程、事件顺序与速度估计会比静态空间判断退化更快；
  \item 场景想象任务会揭示模型是否拥有真正的时空 continuity，而不仅是局部视觉匹配。
\end{itemize}

\subsubsection{错误分析}
建议把错误归因为三层：
\begin{itemize}
  \item \textbf{感知层失败}: 根本没有读到局部几何或事件；
  \item \textbf{结构层失败}: 局部答案互相矛盾，无法形成统一场景；
  \item \textbf{时序层失败}: 知道单帧内容，但无法维持跨时间连续表征。
\end{itemize}

\subsection{第 8 节：讨论与科学解释}
讨论部分建议围绕三句话展开：
\begin{itemize}
  \item \textbf{VLM 看见了什么}: 更偏局部视觉线索、弱物体级几何结构，还是已经形成可重建场景；
  \item \textbf{VLM 没看见什么}: 稳定的参照系、时间连续性、遮挡下的对象持久性、未来状态可预测性；
  \item \textbf{为什么这种分析重要}: 因为它把“模型是否理解时空”从口头判断变成了可重建、可对照、可统计检验的问题。
\end{itemize}

\subsection{第 9 节：局限性与后续工作}
建议坦诚写出：
\begin{itemize}
  \item 当前动态实现更接近 2.5D，而非 full 3D 物理世界；
  \item 目前更强的是行为反演，不是内部神经机制解释；
  \item 更长视频、真实世界视频、人类基线和多摄像机动态场景是下一阶段工作。
\end{itemize}

\section{图表与表格规划}

\subsection{推荐图}
\begin{itemize}
  \item \textbf{Figure 1}: 整体方法图。输入图像/视频 $\rightarrow$ 回答 $\rightarrow$ 约束 $\rightarrow$ 重建 $\rightarrow$ 时空信念分析。
  \item \textbf{Figure 2}: 正确图像、错误图像、无图像三条件下的同一场景重建对比。
  \item \textbf{Figure 3}: 多视角和单视角在 scene belief reconstruction 上的 spread/NRMS 分布。
  \item \textbf{Figure 4}: 动态任务中原序、乱序、反转的性能曲线。
  \item \textbf{Figure 5}: 场景重绘、遮挡补全、未来想象的可视化案例。
\end{itemize}

\subsection{推荐表}
\begin{itemize}
  \item \textbf{Table 1}: 模型总体准确率与缺失率。
  \item \textbf{Table 2}: 重建指标总表（CSR、NRMS、$K_{\text{geom}}$、spread）。
  \item \textbf{Table 3}: 三条件视觉增益与语言先验分析。
  \item \textbf{Table 4}: 时间知觉结果（顺序、速度、轨迹、事件）。
  \item \textbf{Table 5}: 人类与 VLM 在直观时空任务上的差距。
\end{itemize}

\section{最小可发版本建议}

\paragraph{推荐的 v1 论文组合}
\begin{itemize}
  \item \textbf{主线}: 静态空间 benchmark + behavior reconstruction + three-condition visual gain；
  \item \textbf{强补充}: 动态时间知觉小规模主实验（顺序、速度、轨迹）；
  \item \textbf{案例分析}: 场景重绘、遮挡补全、未来想象三类可视化。
\end{itemize}

\paragraph{不建议在 v1 中过度扩展的内容}
\begin{itemize}
  \item full 3D 物理建模；
  \item 过重的内部激活解释；
  \item 过多模型训练细节；
  \item 过宽的真实视频 benchmark 拼盘。
\end{itemize}

\paragraph{一句话结论模板}
即使在 3D 渲染的受控环境和对人类而言高度直观的时空任务中，当前 VLM 也往往只能给出局部正确、全局不稳的回答；只有将其行为转化为可重建的时空信念，我们才能真正回答“模型究竟看见了什么”。

\bibliographystyle{unsrt}
\bibliography{spatiotemporal_vlm_refs}

\end{document}
